\honest{Righting a Wrong Question}{Eliezer Yudkowsky}{2008}

Faced with a question to which you cannot even imagine an answer, try a simple trick. Instead of ``Why do I have free will?'', ask ``Why do I think I have free will?'' The second question is ``guaranteed to have a real answer,'' whether or not there is any such thing as free will. You do believe it, and some chain of causes led to the belief. \nb{Hume gave a similar rule in 1748: when a term is suspected of having no meaning, ``we need but enquire, from what impression is that supposed idea derived?''}

The same works for time, for being born as yourself, for consciousness and for the existence of reality: ask why you think so. \nb{For consciousness, David Chalmers later called this the meta-problem. Whether answering it settles the original question is disputed. Illusionists typically hold that it does; Chalmers holds that explaining our judgments about consciousness ``does not suffice to solve or dissolve the problem of consciousness.''}

I am wearing socks. Tracing back my belief, I go from my visual cortex to my retina, to light reflected from the socks, to the socks, and to why I put them on. My belief is fully explained by the fact that I am wearing socks. A mirage of a lake is explained without any lake: the belief is ``explained away.'' Either way the belief is a real event with a causal history.

You cannot just say ``I believe in free will because I have free will.'' Every step in the sock chain can be explained by biology, optics, electromagnetism or thermodynamics. You must break the chain into small steps, each ``not themselves confusing.'' So where is the free-will sensor in your brain, and how does it work? I leave you with the question.

If your belief comes from observing something real, tracing back will ``eventually'' reach it. If it comes from confusion, you will find an algorithm that runs skew to reality. ``Either way, the question is guaranteed to have an answer.''

Cognitive science may seem less lofty than metaphysics, but at least its questions are ``solvable.'' And the idea that it is less lofty ``is simply wrong.''
